\documentclass[12pt]{article}
\usepackage[utf8]{inputenc}
\usepackage[T1]{fontenc}
\usepackage{hyperref}
\title{Wuzlam Language Resources for K-12 Curriculum Development}
\author{Compiled by Claude Sonnet 5 (claude-sonnet-5) for Andr\'as Kornai}
\date{}
\begin{document}
\maketitle
\section{Introduction}
Wuzlam (ISO 639-3 udl; also Uzlam, Uldeme, Ouldeme, Udlam, Uzam, Mizlime) is a Chadic (Biu-Mandara) language of the Far North region of Cameroon, in the Mandara mountains south of Mora. The most-cited count is about 10,500 speakers (1982). Joshua Project gives 36,000 for the people, which is a different quantity.
\section{Primary Educational Resources}
No school textbooks were found. One source says the language is used in all domains except education, and that literacy classes began in 1999 and produced a dictionary, a grammar and texts.
\section{Linguistic and Reference Materials}
\begin{itemize}
\item Wikipedia article on the Wuzlam language: \url{https://en.wikipedia.org/wiki/Wuzlam_language}
\item Webonary Ouldeme dictionary site, introduction: \url{https://www.webonary.org/ouldeme/overview/introduction/?lang=en}
\item ASJP language page: \url{https://asjp.clld.org/languages/WUZLAM}
\end{itemize}
\section{Religious and Community Resources}
\begin{itemize}
\item Joshua Project profile of the Wuzlam/Uzam: \url{https://joshuaproject.net/people_groups/15428}
\end{itemize}
\section{Digital and Audio Resources}
\begin{itemize}
\item Wikitongues page for udl: \url{https://wikitongues.org/languages/udl/}
\end{itemize}
\section{Curriculum Development Context}
A dictionary, grammar and texts appear to exist from the literacy programme, and the Webonary site is the best lead. Their contents and availability were not examined in this search and should be verified. There is no known formal school use.
\end{document}